%% ALMOST DONE -- edit only where something stated here has actually changed, and then minimally.
%% docs_style.md section 1 is the rule.
\section{Optimization problem}\label{sec:sp:opt}

\subsection{The planner's problem}\label{sec:sp:opt:problem}

\smalltitle{Controls and states.} The planner chooses paths for the eight controls
\begin{align}
    \left\{C,\;I,\;N,\;D,\;\varpi,\;K^R,\;W,\;\Omega\right\},
\label{eq:sp:controls}
\end{align}
which determine the remaining flows through the definitions
\begin{align}
H &= \mu\mathcal W,
&
R^R &= \mathcal R\left(\varpi H,K^R, t\right),
&
R &= N+R^R,
&
Y &= F\left(K-K^R,\,R,\,P, t\right),
&
G &= I+\delta K,
\label{eq:sp:definitions}
\end{align}
and it carries the six states
\begin{align}
\left\{K,\;S,\;X,\;P,\;M^K,\;\mathcal W\right\},
\label{eq:sp:states}
\end{align}
with $K_0,S_0,X_0,P_0,M^K_0,\mathcal W_0$ given.

Two of the eight controls are not free. Given the remaining six, the ledger~\eqref{eq:sp:ledger} and the intensity condition~\eqref{eq:sp:intensity} are two equations in $W$ and $\Omega$, and the goods constraint~\eqref{eq:sp:goods} removes one more degree of freedom, leaving five independent instruments against the six states. The waste stock is the state without an instrument of its own: its inflow is whatever generation the other choices imply, and its outflow is mechanical at the handling share $\mu$. We nonetheless carry $W$ and $\Omega$ as controls with the two relations adjoined, so that their multipliers $p^W$ and $\zeta$ appear explicitly and can be read as shadow prices.


\begin{problem}[The planner's problem]\label{prob:sp}
Choose $\left\{C_t,I_t,N_t,D_t,\varpi_t,K^R_t,W_t,\Omega_t\right\}_{t\ge0}$ to maximise

\noindent $U_0=\sum_{t=0}^\infty\beta^t\left[u(C_t)-v(P_t)\right]$ subject to, in every period $t$,
\begin{align}
&\text{\emph{goods:}}
&
F\left(K-K^R,R,P,t\right) &= C+I+\delta K+C^N(N,S,t)+C^D(D,X,t)+c^W\!\left(\varpi,t\right)H,
\label{eq:sp:goods}
\\[0.2em]
&\text{\emph{ledger:}}
&
W &= R-\Omega\phi^I G+\delta M^K+\mathcal N,
\label{eq:sp:ledger}
\\[0.2em]
&\text{\emph{intensity:}}
&
R &= \Omega\left(\mathcal D+\phi^IG\right),
\label{eq:sp:intensity}
\end{align}
and to the transition equations
\begin{align}
K_{t+1} &= K_t+I_t,
&
S_{t+1} &= S_t+D_t-N_t,
&
X_{t+1} &= X_t+D_t,
\notag\\
P_{t+1} &= P_t+\Xi_t-\theta(P_t)P_t,
&
M^K_{t+1} &= \left(1-\delta\right)M^K_t+\Omega_t\phi^I G_t,
&
\mathcal W_{t+1} &= \left(1-\mu\right)\mathcal W_t+W_t,
\label{eq:sp:states-motion}
\end{align}
with waste bases $\mathcal{N}$ and $\mathcal{D}$ defined in \eqref{eq:sp:setup:base}, $c^W\!\left(\varpi,t\right) \equiv c^c(t)\varpi+ c^T(\varpi,t)$, the definitions in~\eqref{eq:sp:definitions}, the emission flow from \eqref{eq:sp:setup:Xi}, levied on the handled flow $H$, and domains $C,I+\delta K,N,D,K^R,W,\Omega\ge0$, $\varpi\in\left[0,1\right]$.
\end{problem}

The two static relations adjoined alongside the goods constraint are the accounting identities of Section~\ref{sec:sp:setup:aggregatematerial}, restated in the variables the planner chooses. Equation~\eqref{eq:sp:ledger} is the ledger~\eqref{eq:sp:setup:ledger}: waste generated equals the material entering production, less the net accumulation of material embodied in the capital stock, plus the mass displaced from nature. Equation~\eqref{eq:sp:intensity} is the balance of embodied materials in the final good, \eqref{eq:sp:setup:ybalance_OmegaPhi}, and is what determines $\Omega$. The two are distinct restrictions rather than two forms of one identity: substituting the second into the first returns
\begin{align}
W &= \Omega\mathcal D+\delta M^K+\mathcal N,
\label{eq:sp:ledger-streams}
\end{align}
which is the sum of the individual waste streams~\eqref{eq:sp:setup:totalwaste_generic}. The pair thus carries no content beyond the process-level accounting, but together they pin down $W$ and $\Omega$, as noted above.

$\Omega$ itself has a somewhat complicated interpretation, as it captures the level of the average material intensity of production while also adjusting to reflect whether final goods are put to more or less material-intensive purposes. From~\eqref{eq:sp:intensity}, $\Omega=R/\left(\mathcal D+\phi^IG\right)$ scales with $R$ for a fixed structure of uses, but for a given $R$ it falls whenever $Y$ is shifted towards a use carrying a larger coefficient ($\omega^j$ or $\phi^I$). We refer to the relative sizes of the different uses of the final good as the \textit{composition of activities/uses} of the good.

The channel through which that composition reaches the allocation is narrow, and worth marking before the optimality conditions. Given $R$, the ledger says that everything entering production becomes waste except what gross investment stores in $M^K$. The composition therefore matters only by moving $\Omega$ through $\Phi^I=\Omega\phi^I$ i.e. the quantity of material held back from today's waste flow. This is precisely what the multiplier $\zeta$ on~\eqref{eq:sp:intensity} prices, and Section~\ref{sec:sp:opt:foc} shows that the channel closes entirely when $\phi^I=0$.

Two features of the primitives of Section~\ref{sec:sp:setup:primitives} deserve a word before the conditions are derived. First, exogenous technological change makes the problem non-autonomous: $F$, $\mathcal R$, $C^N$, $C^D$ and $c^W$ all carry the time argument $t$. Since $t$ is not a choice variable, it generates no first-order condition of its own; every condition below holds period by period with the primitives evaluated at that date, and following the convention of Section~\ref{sec:sp:setup:primitives} we suppress the argument from here on. Second, the exhaustibility of discoveries, \eqref{eq:sp:setup:exhaustible-discovery}, never has to be imposed on the problem as a constraint $X\le\bar X_{\max}$: as $X$ approaches the ceiling, the marginal cost of even the first tonne of discovery diverges, so the bound enforces itself through the exploration margin. This is also why the exploration condition of Section~\ref{sec:sp:opt:foc} is stated as a corner--interior system from the outset: the corner is not a boundary case to be checked but a phase of every solution, and the phase structure it induces is recorded there.


Next, we note that $\Xi$ is written as linear in $H$ and $R^R$ separately, so that no derivative of $a(\cdot)$ ever appears inside a pollution term. All the curvature of the recycling
technology enters through $\mathcal R$ and through $\alpha$ and $a'$ alone.\footnote{Recall
from~\eqref{eq:sp:setup:recycling-derivatives} that $\mathcal R_T=\alpha\equiv a-a'x$ is the marginal
recycling yield and $\mathcal R_{K^R}=a'$.}


\smalltitle{The Lagrangian.} Write the discounted Lagrangian with the static constraints
adjoined. It is convenient to normalise every multiplier by $\lambda_t$, the multiplier on the goods
constraint~\eqref{eq:sp:goods}, so that all shadow prices below are measured in units of the final
good rather than in utils. Each transition equation in~\eqref{eq:sp:states-motion} is adjoined in the period in which its choices are made, so the costate dated $t$ prices the stock those choices determine --- the beginning-of-period-$t{+}1$ stock. Accordingly, let the costates be
\begin{align}
\underbrace{\lambda q}_{K},
&&
\underbrace{\lambda p^S}_{S},
&&
\underbrace{\lambda p^X}_{X},
&&
\underbrace{-\lambda p^P}_{P},
&&
\underbrace{-\lambda p^M}_{M^K},
&&
\underbrace{\lambda p^{\mathcal W}}_{\mathcal W},
\label{eq:sp:costates-def}
\end{align}
the two minus signs anticipating that both $P$ and $M^K$ are likely liabilities, so that $p^P$ and $p^M$ read as shadow \emph{costs}; the stockpile costate carries a plus sign, anticipating that on the paths of interest the anthropogenic reserve is an \emph{asset}, though nothing rules out $p^{\mathcal W}<0$ when handling is a net burden. Let $\lambda p^W$ be the multiplier on the ledger~\eqref{eq:sp:ledger} written with $W$ on the left, and $\lambda\zeta$ the multiplier on the intensity definition~\eqref{eq:sp:intensity} written with
$R$ on the left. Then
\begin{align}
\mathcal L
={}& \sum_{t=0}^{\infty}\beta^t\bigg\{ u(C_t)-v(P_t)
\notag\\
&+\lambda_t\Big[F\left(K_t-K^R_t,R_t,P_t,t\right)-C_t-I_t-\delta K_t-C^N(N_t,S_t,t)-C^D(D_t,X_t,t)-c^W\!\left(\varpi_t,t\right)\mu\mathcal W_t\Big]
\notag\\
&+\lambda_t q_t\left[K_t+I_t-K_{t+1}\right]
\;+\;\lambda_t p^S_t\left[S_t+D_t-N_t-S_{t+1}\right]
\;+\;\lambda_t p^X_t\left[X_t+D_t-X_{t+1}\right]
\notag\\
&-\lambda_t p^P_t\left[P_t+\Xi_t-\theta(P_t)P_t-P_{t+1}\right]
\;-\;\lambda_t p^M_t\left[\left(1-\delta\right)M^K_t+\Omega_t\phi^IG_t-M^K_{t+1}\right]
\notag\\
&+\lambda_t p^{\mathcal W}_t\left[\left(1-\mu\right)\mathcal W_t+W_t-\mathcal W_{t+1}\right]
\;+\;\lambda_t p^W_t\Big[W_t-N_t-R^R_t+\Omega_t\phi^IG_t-\delta M^K_t-\mathcal N_t\Big]
\notag\\
&+\lambda_t \zeta_t\Big[N_t+R^R_t-\Omega_t\left(\mathcal D_t+\phi^IG_t\right)\Big]\bigg\},
\label{eq:sp:lagrangian}
\end{align}
with relevant definitions substituted throughout.\footnote{This includes $R=N+R^R$, $H=\mu\mathcal W$, $R^R=\mathcal R\left(\varpi H,K^R,t\right)$, $G=I+\delta K$, $Y=F\left(K-K^R,R,P,t\right)$
and $C^W=c^W\!\left(\varpi,t\right)H$ substituted throughout.} Here, the multiplier $p^W$ measures the \textit{social cost} of waste in units of final goods and $\zeta$ measures the social value of embodied materials in units of final goods.\footnote{To see this, assume that we could increase waste from nature $\mathcal{N}$ marginally without changing any other controls, i.e. if more waste was generated out of the blue; we would have $\partial \mathcal L^*/\partial \mathcal N = - \lambda p^W$ which would be negative (i.e. a cost) exactly if $p^W>0$. Similarly, perturb the condition $R + \kappa = \Omega \left(\mathcal{D}+\phi^IG\right)$ such that we could increase embodied materials (right hand side) marginally without increasing the amount of materials $R$ that goes into the production process; we would then have $\partial \mathcal L^*/\partial \kappa = \lambda \zeta$ which would be positive if $\zeta>0$.}

The dating convention in~\eqref{eq:sp:costates-def} does real work in one place: a tonne of waste generated in period $t$ enters the stock at $t+1$, and $p^{\mathcal W}_t$ is by construction the period-$t$ value of exactly such a tonne, which is what keeps the waste margin below a contemporaneous identification rather than a discounted one.



\subsection{Optimality conditions}\label{sec:sp:opt:foc}

The conditions are collected here; Appendix~\ref{sec:app:sp} derives each of them, decomposes it into the separate margins it balances, and checks the units. Define
\begin{align}
\Psi &\equiv F_R\left(1-\zeta\Omega^Y\right)+\zeta,
&
\sigma &\equiv \frac{\Phi^IG}{R}\;\in\left[0,1\right],
&
1+r_{t+1} &\equiv \frac{\lambda_t}{\beta\lambda_{t+1}},
\label{eq:sp:shorthand}
\end{align}
where $\Psi$ is the value of a tonne of raw material delivered to production, $\sigma$ is the share of material embodied in new capital, and $r_{t+1}$ is the social discount rate between periods $t$ and $t+1$ implied by the goods multiplier (so $r=\rho$ whenever $\lambda$ is constant). $\Psi$ captures output net of embodied material output requires plus relaxation of the intensity condition, but excludes the ledger: the tonne must eventually leave as waste, and that is priced separately by $p^W$. With a minimum material requirement, $F_R$ means $\mathcal F_{R^e}$ evaluated at the prevailing $R\ge\bar R$. The region $R<\bar R$ produces nothing and is never operated --- paying for material without clearing the floor is strictly dominated --- so no condition below is ever evaluated there; the genuine alternative that region represents, not operating at all, is a discrete comparison rather than a first-order event, and Section~\ref{sec:sp:opt:foc} returns to it under ``the material margin''.

\smalltitle{Controls.} Each condition below is a within-period condition, holding in every period $t$ with all quantities dated $t$: every control acts on flows and constraints of its own period alone --- the one intertemporal consequence a control has, the stock its choice feeds, is priced by a costate dated the same period --- so no discount factor or lead appears in any of them, and only the costate equations further below carry the passage of time. Each condition is stated with the margin it balances. Appendix~\ref{sec:app:sp} derives each one, decomposes it term by term, and checks the units.

\emph{Material intensity.} Raising $\Omega$ stores material in the capital stock rather than releasing it now. The gain is the price gap $p^W-p^M$, scaled by the share $\sigma$ of material that storage actually captures, and the shadow price of embodied material must equal it:
\begin{align}
\zeta &= \sigma\left(p^W-p^M\right).
\label{eq:sp:sum:Omega}
\end{align}

\emph{Consumption.} A unit consumed absorbs one final good and carries $\Omega^C$ tonnes of embodied material away from the storage margin. Marginal utility therefore exceeds $\lambda$ by that wedge, so $\lambda$ is not the marginal utility of consumption in this model:
\begin{align}
u'\left(C\right)/\lambda &= 1+\zeta\Omega^C .
\label{eq:sp:sum:C}
\end{align}

\emph{Investment.} A unit invested buys a unit of capital and stores $\Phi^I$ tonnes, of which the fraction $\sigma$ is clawed back because embodying material also consumes material drawn through $R$. The shadow price of installed capital departs from unity by exactly the net storage value:
\begin{align}
q &= 1-\Phi^I\left(1-\sigma\right)\left(p^W-p^M\right).
\label{eq:sp:sum:I}
\end{align}

\emph{Waste.} In the period of generation a tonne of waste has exactly one consequence: it enters next period's stockpile. Collection, treatment, recycling and emissions are all levied on the handled flow $H=\mu\mathcal W$ and are therefore priced on the state, not on the control, so the waste margin collapses to an identification:
\begin{align}
p^W_t &= -p^{\mathcal W}_t .
\label{eq:sp:sum:W}
\end{align}
The social cost of generating a tonne of waste is minus the value of a tonne in the anthropogenic reserve, and under the dating convention of~\eqref{eq:sp:costates-def} the identification is exact and contemporaneous: $p^{\mathcal W}_t$ is the period-$t$ price of a tonne that sits in the stock from $t+1$ onward, which is precisely what a tonne generated in period $t$ is. Everything levied on handling --- resources absorbed, emissions, feedstock recovered --- sits in the stockpile costate~\eqref{eq:sp:sum:Wstock} below, and waste generation is a social cost precisely when the stockpile is a liability.

\smalltitle{The four bounded controls.} The four conditions above are stated as equalities, and for
those controls this is without loss at any candidate optimum: consumption because we maintain the
Inada condition $u'\left(C\right)\to\infty$ as $C\to0$, investment and the two accounting controls
because their bounds are never approached. The remaining four controls are not in that class.
Extraction and exploration are bounded below, $N\ge0$ and $D\ge0$, and each corner is reachable;
under exhaustible discoveries~\eqref{eq:sp:setup:exhaustible-discovery} the exploration corner is
not merely reachable but eventually unavoidable, so the $D=0$ regime is a phase of the solution
rather than a boundary curiosity. Treatment is bounded on both sides, $\varpi\in\left[0,1\right]$,
and recycling capital is bounded below, $K^R\ge0$; both bounds are reachable. All four conditions
are therefore stated as Kuhn--Tucker systems: each corner is reported together with the inequality
that makes it optimal, and the interior equality together with the inequality that makes \emph{it}
optimal, so that the branches partition the possibilities. Finally, the state bound $\mathcal W\ge0$
never needs a multiplier: with $W\ge0$ and $\mu<1$ the stock obeys $\mathcal W_t\ge\left(1-\mu\right)^t\mathcal W_0>0$,
so the stockpile cannot empty in finite time and the recycling block never runs out of feedstock.

\emph{Extraction.} A tonne extracted is worth $\Psi$ on delivery. It must cover the extraction
effort, grossed up for the embodied material the effort itself carries, the in-situ rent, and the
waste it becomes (both the tonne and the overburden displaced alongside it):
\begin{subequations}
\label{eq:sp:sum:Nsys}
\begin{align}
N &= 0,
\qquad \text{if}\quad
\Psi \;\le\; C^N_N\left(0,S,t\right)\left(1+\zeta\Omega^N\right)+p^S+p^W\left(1+\Omega^{N,S}\right),
\label{eq:sp:sum:N-corner}
\\
\Psi &= C^N_N\left(1+\zeta\Omega^N\right)+p^S+p^W\left(1+\Omega^{N,S}\right),
\qquad \text{otherwise.}
\label{eq:sp:sum:N}
\end{align}
\end{subequations}
The corner is not excluded by an Inada condition on materials, as it would be in a model without
recycling: $F_R\to\infty$ as $R\to\bar R^{+}$ puts an unbounded value on the marginal tonne of $R$
at the threshold, but not on the marginal \emph{virgin} tonne, because $R^R>0$ may keep $R$ away
from that limit and $\Psi$ finite. Note that with $\bar R>0$ the Inada condition, when imposed at
all, bites at a strictly positive material flow rather than at the origin, so it no longer implies
that material use must remain positive forever.
Against that finite benefit, the right-hand side of~\eqref{eq:sp:sum:N-corner} carries the strictly
positive floor $p^W\left(1+\Omega^{N,S}\right)$: the waste charge on the tonne and its overburden is
levied per tonne extracted and does not vanish as $N\to0$. Whether the corner is in fact reached is
a question about the essentiality of materials, taken up in Section~\ref{sec:sp:longrun}.

\emph{Exploration.} A tonne discovered creates the rent $p^S$ and uses up a deposit, $p^X<0$; the
net value must cover exploration effort, grossed up for the embodied material the effort carries,
and the waste charge on the mass exploration displaces:
\begin{subequations}
\label{eq:sp:sum:Dsys}
\begin{align}
D &= 0,
\qquad \text{if}\quad
p^S+p^X \;\le\; C^D_D\left(0,X,t\right)\left(1+\zeta\Omega^D\right)+p^W\Omega^{D,S},
\label{eq:sp:sum:D-corner}
\\
p^S+p^X &= C^D_D\left(1+\zeta\Omega^D\right)+p^W\Omega^{D,S},
\qquad \text{otherwise.}
\label{eq:sp:sum:D}
\end{align}
\end{subequations}
This is the corner the primitives force. Under~\eqref{eq:sp:setup:exhaustible-discovery} the
right-hand side of~\eqref{eq:sp:sum:D-corner} diverges as $X\to\bar X_{\max}$ while the left-hand
side stays bounded ($p^X\le0$, and $p^S$ is finite along any admissible path; both are priced among
the costates below), so exploration must die out --- either $D$ hits the corner outright or dwindles
as $X$ approaches its ceiling --- and $X$ never in fact reaches $\bar X_{\max}$.
Appendix~\ref{sec:app:sp} records what the cessation of exploration does to the discovery costate.

\emph{The recycling block.} To keep the branches of the remaining two systems readable, write the
marginal gain on each of the two margins as
\begin{align}
\mathcal G^{\varpi} &\equiv \alpha\left(\Psi-p^W\right)+p^P\left[\left(1-d^W\right)+d^W\alpha\right],
&
\mathcal G^{K} &\equiv \Psi-p^W+d^Wp^P,
\label{eq:sp:sum:gains}
\end{align}
both in final goods per tonne: $\mathcal G^\varpi$ is the value of routing one more handled tonne to
treatment, and $\mathcal G^{K}$ the value of one more tonne of secondary material recovered by
capital, so that the gain from a unit of recycling capital is $a'\!\left(x\right)\mathcal G^K$.

\emph{Treatment share.} Treating one more tonne yields $\alpha$ tonnes of secondary material, each
worth $\Psi-p^W$ rather than $\Psi$ because the recovered material must itself leave as waste later,
and it diverts the tonne from the environment. The gain $\mathcal G^\varpi$ is set against a rising
marginal handling cost $c^c+c^{T\prime}\left(\varpi\right)$, and since $c^{T\prime}\left(0\right)=0$
by~\eqref{eq:sp:setup:waste-cost} the marginal cost of the first treated tonne is the collection
charge alone, making the lower corner a genuine choke at $c^c$:
\begin{subequations}
\label{eq:sp:sum:varpi}
\begin{align}
\varpi &= 0,
&&\text{if}\quad \mathcal G^{\varpi} \;\le\; c^c\left(1+\zeta\Omega^W\right),
\label{eq:sp:sum:varpi-lower}
\\
\mathcal G^{\varpi} &= \left[c^c+c^{T\prime}\left(\varpi\right)\right]\left(1+\zeta\Omega^W\right),
&&\text{if}\quad c^c\left(1+\zeta\Omega^W\right) \;<\; \mathcal G^{\varpi} \;<\; \left[c^c+c^{T\prime}\left(1\right)\right]\left(1+\zeta\Omega^W\right),
\label{eq:sp:sum:varpi-interior}
\\
\varpi &= 1,
&&\text{if}\quad \mathcal G^{\varpi} \;\ge\; \left[c^c+c^{T\prime}\left(1\right)\right]\left(1+\zeta\Omega^W\right).
\label{eq:sp:sum:varpi-upper}
\end{align}
\end{subequations}
Convexity of $c^T$ makes the marginal cost strictly increasing, so the three branches are exhaustive
and mutually exclusive, and the interior root is unique when it exists. The lower corner is
economically meaningful --- the economy dumps its handled waste untreated whenever the gain cannot
cover collection --- and charging collection per treated tonne, with dumping free, is what keeps that
corner feasible for an economy whose output has collapsed. The upper corner is reachable
by assumption: $c^{T\prime}\left(1\right)<\infty$ is maintained throughout
by~\eqref{eq:sp:setup:cT-finite}, so the third branch is non-empty at every date and full treatment
occurs whenever the gain is large enough. This is deliberate. Since only the product $a\varpi$ enters
the ledger, a treatment cost diverging at $\varpi=1$ would cap circularity on its own and would
silently do the work of a hard yield ceiling; ruling it out leaves the recycling technology as the
only source of a circularity limit.

\emph{Recycling capital.} A final good moved into recycling yields $a'$ tonnes, valued at $\mathcal G^K$. It costs the marginal product forgone in production, net
of the material that output would have required. The Kuhn-Tucker conditions then read:
\begin{subequations}
\label{eq:sp:sum:KR}
\begin{align}
K^R &= 0,
&&\text{if}\quad a'\!\left(0\right)\mathcal G^{K} \;\le\; F_K\left(1-\zeta\Omega^Y\right),
\label{eq:sp:sum:KR-corner}
\\
a'\!\left(x\right)\mathcal G^{K} &= F_K\left(1-\zeta\Omega^Y\right),
&&\text{if}\quad a'\!\left(0\right)\mathcal G^{K} \;>\; F_K\left(1-\zeta\Omega^Y\right).
\label{eq:sp:sum:KR-interior}
\end{align}
\end{subequations}
The upper bound $K^R\le K$ is
never stated because it never binds: under the maintained Inada condition $F_K\to\infty$ as
$K^Y\to0$, the corner
violates~\eqref{eq:sp:sum:KR-interior} long before it is reached.

\emph{The two recycling conditions restated.} Whenever extraction is interior, the combination
$\Psi-p^W$ inside $\mathcal G^\varpi$ and $\mathcal G^K$ can be eliminated using the extraction
condition~\eqref{eq:sp:sum:N}, and the unit waste charge cancels (see appendix \ref{sec:app:sp}). Define the \emph{virgin displacement value}
\begin{align}
V &\equiv \Psi-p^W+d^Wp^P \;=\; C^N_N\left(1+\zeta\Omega^N\right)+p^S+\Omega^{N,S}p^W+d^Wp^P,
\label{eq:sp:sum:V}
\end{align}
we have $\mathcal G^K=V$ and $\mathcal G^\varpi=\alpha V+p^P\left(1-d^W\right)$, so that the two
Kuhn--Tucker systems read
\begin{subequations}
\label{eq:sp:sum:recycling-virgin}
\begin{align}
\varpi &= 0
&&\text{if}\quad \alpha V+p^P\left(1-d^W\right)\le c^c\left(1+\zeta\Omega^W\right),
\notag
\\
\alpha V+p^P\left(1-d^W\right) &= \left[c^c+c^{T\prime}\left(\varpi\right)\right]\left(1+\zeta\Omega^W\right)
&&\text{if}\quad c^c\left(1+\zeta\Omega^W\right)<\alpha V+p^P\left(1-d^W\right)<\left[c^c+c^{T\prime}\left(1\right)\right]\left(1+\zeta\Omega^W\right),
\notag
\\
\varpi &= 1
&&\text{if}\quad \alpha V+p^P\left(1-d^W\right)\ge \left[c^c+c^{T\prime}\left(1\right)\right]\left(1+\zeta\Omega^W\right),
\label{eq:sp:sum:recycling-virgin-varpi}
\end{align}
and
\begin{align}
K^R &= 0
&&\text{if}\quad a'\!\left(0\right)V\le F_K\left(1-\zeta\Omega^Y\right),
\notag
\\
a'\!\left(x\right)V &= F_K\left(1-\zeta\Omega^Y\right)
&&\text{if}\quad a'\!\left(0\right)V> F_K\left(1-\zeta\Omega^Y\right).
\label{eq:sp:sum:recycling-virgin-KR}
\end{align}
\end{subequations}
\emph{Recycling is paid the cost of the virgin tonne it displaces, and nothing else}: extraction
effort with the material it carries, the in-situ rent, the overburden that tonne would have brought
up, and the leakage it would have suffered. The charge $p^W$ is absent because the recovered tonne and
the virgin tonne it replaces face the same charge on the way out; recycling changes the route a
tonne takes into production, never the fact that it must eventually leave. The one motive $V$ cannot
price is the diversion of the untreated tonne from the environment, $p^P\left(1-d^W\right)$, which is
why that term survives separately on the treatment margin and has no counterpart on the capital
margin. The restatement is phase-specific: in a phase with $N=0$, condition~\eqref{eq:sp:sum:N}
holds only as an inequality, the elimination behind $V$ is unavailable, and the primitive
forms~\eqref{eq:sp:sum:varpi} and~\eqref{eq:sp:sum:KR} are the ones to carry.

\smalltitle{The material margin.} With $\bar{R}>0$, the above conditions only hold for an economy without material collapse (i.e. without $R<\bar{R}$ and thus $Y= 0$). We do not further analyse paths that do end up in material collapse, our analysis only points to conditions where such collapse is inevitable.


\smalltitle{Phase structure.} The two corners outside the recycling block do not have the same
status, and the difference organises the optimal path into material regimes. The exploration corner
is forced by the primitives, as recorded above, so the $D=0$ regime must be carried as a genuine
phase of the solution. The extraction corner is a possibility, not a necessity: whether $N=0$ is
reached depends on whether the value of the marginal tonne stays finite where the corner
inequality~\eqref{eq:sp:sum:N-corner} could bind, and Section~\ref{sec:sp:longrun} shows this to be
exactly the choice between ending the material era by \emph{stranding} (the corner binds in finite
time, $S_\infty>0$) and ending it by \emph{exhaustion} ($N\to0$ only asymptotically, $S\to0$).
Either way, the demise of exploration is what gives the model its long-run bite: cumulative virgin
extraction is bounded by the material that can ever be known,
\begin{align}
\sum_{t=0}^{\infty} N_t \;\le\; S_0+\bar X_{\max}-X_0,
\label{eq:sp:sum:cumulative-N}
\end{align}
so virgin input cannot remain bounded away from zero forever, and whatever raw material the economy
uses in the long run must come from recycling. With $\bar R>0$ a third route to the end of the material era opens alongside stranding and exhaustion, and it is neither: the discrete shutdown in which material use stops at a strictly positive flow because the remaining stretch of the technology is dominated by not operating at all. Under the hard reading of the floor this route is drastic: permanent shutdown ends goods production itself, so it is taken
only when holding the floor has become infeasible.



\smalltitle{Costates.} Each state $Z$ carries the costate $m$ named in~\eqref{eq:sp:costates-def}, with the transversality condition $\lim_{t\to\infty}\beta^t\lambda_t m_t Z_{t+1}=0$. That condition is what selects the forward solution reported alongside each recursion below, and it gives every price the same asset-pricing form: the discounted stream of the dividend the stock pays while it is held. Every recursion has the same shape --- $\left(1+r_{t+1}\right)m_t$ equals next period's dividend plus the surviving fraction of the stock times next period's price --- and we introduce the compound discount factor
\begin{align}
\Lambda_{t,s} &\equiv \prod_{j=t+1}^{s}\frac{1}{1+r_j},
\label{eq:sp:sum:discount}
\end{align}
so that every forward solution below is a $\Lambda$-weighted sum of dividends.

\emph{Capital.} A unit of capital pays its marginal product, net of the embodied material the extra output requires, and survives to the next period at rate $1-\delta$:
\begin{align}
\left(1+r_{t+1}\right)q_t &= F_K\left(1-\zeta\Omega^Y\right)\Big|_{t+1}+\left(1-\delta\right)q_{t+1},
\qquad
q_t=\sum_{s=t+1}^\infty \Lambda_{t,s}\left(1-\delta\right)^{s-t-1}F_K\left(1-\zeta\Omega^Y\right)\Big|_s .
\label{eq:sp:sum:K}
\end{align}
Dividing the recursion by $q_t$ returns the familiar arbitrage form $1+r_{t+1}=\left[F_K\left(1-\zeta\Omega^Y\right)\big|_{t+1}+\left(1-\delta\right)q_{t+1}\right]/q_t$.\footnote{The storage motive of~\eqref{eq:sp:sum:I} appears nowhere in it: depreciation destroys capital and releases its embodied material at the same rate, so the two cancel in the return.}

\emph{Reserves.} A tonne left in the ground pays a dividend by making the extraction still to come cheaper, since $C^N_S<0$:
\begin{align}
\left(1+r_{t+1}\right)p^S_t &= \left|C^N_S\right|\left(1+\zeta\Omega^N\right)\Big|_{t+1}+p^S_{t+1},
\qquad
p^S_t=\sum_{s=t+1}^\infty \Lambda_{t,s}\left|C^N_S\right|\left(1+\zeta\Omega^N\right)\Big|_s .
\label{eq:sp:sum:S}
\end{align}
This is Hotelling's rule with a stock effect: the rent is positive and need not grow at the rate of interest, because the reserve yields a flow benefit as well as a capital gain.

\emph{Discoveries.} Cumulative discovery is a pure liability, raising the cost of all future exploration:
\begin{align}
\left(1+r_{t+1}\right)p^X_t &= -C^D_X\left(1+\zeta\Omega^D\right)\Big|_{t+1}+p^X_{t+1},
\qquad
p^X_t=-\sum_{s=t+1}^\infty \Lambda_{t,s}\,C^D_X\left(1+\zeta\Omega^D\right)\Big|_s .
\label{eq:sp:sum:X}
\end{align}
Its negative sign is what stops exploration from being a costless substitute for extraction in~\eqref{eq:sp:sum:D}. Under exhaustible discoveries the liability expires with the activity it prices: once exploration has ceased for good, $C^D_X$ is evaluated at $D=0$, where $C^D\left(0,X,t\right)=0$ implies $C^D_X\left(0,X,t\right)=0$, so $p^X=0$ from that date onward (Appendix~\ref{sec:app:sp}).

\emph{Pollution.} The stock damages utility and output, the two together giving a marginal damage $\mathcal V\equiv v'\left(P\right)/\lambda-F_P\left(1-\zeta\Omega^Y\right)>0$ in final goods per tonne, and it survives each period at the \emph{marginal} retention factor $1-\theta-\theta'P$ rather than the average factor $1-\theta$:
\begin{align}
\left(1+r_{t+1}\right)p^P_t &= \mathcal V_{t+1}+\left[1-\theta-\theta'P\right]_{t+1}p^P_{t+1},
\qquad
p^P_t=\sum_{s=t+1}^\infty \Lambda_{t,s}\left[\prod_{j=t+1}^{s-1}\left(1-\theta-\theta'P\right)_j\right]\mathcal V_s .
\label{eq:sp:sum:P}
\end{align}
Both damages are costs, so $p^P>0$ unconditionally. Since $\theta'<0$ the marginal decay rate is below the average, so the retention factor is the larger of the two, and if decay saturates far enough that $\left(1-\theta-\theta'P\right)/\left(1+r\right)\ge1$ --- equivalently $r+\theta+\theta'P\le0$ --- the sum need not converge (this model's version of a tipping point).

\emph{Embodied materials.} A tonne held in the capital stock is released at rate $\delta$, becoming waste when it is:
\begin{align}
\left(1+r_{t+1}\right)p^M_t &= \delta\,p^W_{t+1}+\left(1-\delta\right)p^M_{t+1},
\qquad
p^M_t=\sum_{s=t+1}^\infty \Lambda_{t,s}\left(1-\delta\right)^{s-t-1}\delta\,p^W_s .
\label{eq:sp:sum:M}
\end{align}
The integrated form is the one to keep in view: $p^M$ is a \emph{forward} average of $p^W$, not a function of its current value.

\emph{The anthropogenic reserve.} A tonne in the stockpile is drawn for handling with share $\mu$ each period, and each handled tonne yields the \emph{handling dividend}
\begin{align}
h &\equiv \alpha\varpi\left(\Psi-p^W\right)-c^W\left(1+\zeta\Omega^W\right)-p^P\left[1-\varpi\left(1-d^W\right)-d^W\alpha\varpi\right],
\label{eq:sp:sum:h}
\end{align}
in final goods per tonne: the feedstock recovered, net of the waste it will regenerate, less the handling resources and the emission damage. The costate obeys
\begin{align}
\left(1+r_{t+1}\right)p^{\mathcal W}_t &= \mu\,h_{t+1}+\left(1-\mu\right)p^{\mathcal W}_{t+1},
\qquad
p^{\mathcal W}_t=\sum_{s=t+1}^\infty \Lambda_{t,s}\left(1-\mu\right)^{s-t-1}\mu\,h_s .
\label{eq:sp:sum:Wstock}
\end{align}
The stockpile is priced exactly as the virgin reserve is: $S$ pays its dividend by making future extraction cheaper, $\mathcal W$ by making future recovery possible, and both prices are forward sums of those dividends, weighted by the geometric probability that the tonne is still in the stock when the dividend is paid. Two readings of~\eqref{eq:sp:sum:Wstock} are worth recording. First, substituting $p^W=-p^{\mathcal W}$ into $h$ and collecting terms shows that the effective survival factor on the stockpile is $1-\mu\left(1-\alpha\varpi\right)$ rather than $1-\mu$: in marginal-value terms the stockpile runs off only through genuine leakage, because the recovered fraction $\alpha\varpi$ of each handled tonne re-enters production and eventually returns to the stockpile as waste. Second, at a rest point of the price system, \eqref{eq:sp:sum:Wstock} gives $p^{\mathcal W}=\mu h/\left(r+\mu\right)$, and combining with~\eqref{eq:sp:sum:W} yields the rest-point waste margin
\begin{align}
p^W\left(1-\alpha\varpi+\frac{r}{\mu}\right) &= c^W\left(1+\zeta\Omega^W\right)+p^P\left[1-\varpi\left(1-d^W\right)-d^W\alpha\varpi\right]-\alpha\varpi\,\Psi .
\label{eq:sp:sum:W-limit}
\end{align}
Up to the correction $r/\mu$, this is the condition a model without a waste stock would impose in every period; the correction prices the fact that a stockpiled tonne turns over one period at a time, so that the margins levied on it are paid with delay and discounting. The waste stock changes \emph{when} the waste price is paid, and at a rest point that timing survives only as the $r/\mu$ term.\footnote{At the buffer corner $\mu=1$ the recursion collapses to $p^W_t=-h_{t+1}/\left(1+r_{t+1}\right)$: the entire static waste margin, evaluated one period ahead and discounted --- waste generated today is handled in full tomorrow, and priced accordingly.}

\emph{Consumption.} Combining~\eqref{eq:sp:sum:C} across two adjacent periods with the definition of $r$ gives the Euler equation with one extra factor:
\begin{align}
\frac{u'\left(C_t\right)}{\beta\,u'\left(C_{t+1}\right)} &= \left(1+r_{t+1}\right)\frac{1+\zeta_t\Omega^C_t}{1+\zeta_{t+1}\Omega^C_{t+1}},
\qquad\text{i.e.}\qquad
\eta\,\Delta\ln C_{t+1} = \ln\frac{1+r_{t+1}}{1+\rho}-\Delta\ln\left(1+\zeta\Omega^C\right)_{t+1}
\label{eq:sp:sum:Cgrowth}
\end{align}
under CRRA utility with $\eta\equiv-u''C/u'$. A consumption wedge that is constant over time affects only the level of consumption; it is a wedge that is \emph{changing} that tilts the path.



\smalltitle{Reading the system.} Three prices carry the materials side of the problem, and every condition above is built from them together with the standard terms. $\Psi$ is what a tonne is worth entering production; $p^W$ is what it costs on the way out; and $\zeta$ is what it is worth to have a tonne stored in the capital stock rather than diffused. The separation is deliberate. A tonne is valued once when it arrives, charged once when it leaves, and the difference between the two dates is priced by $\zeta$. With the waste stock, ``the way out'' is itself an asset: $p^W=-p^{\mathcal W}$, so the exit charge is negative exactly when the stockpile the tonne joins is valuable.

Relative to the same model without materials accounting, two patterns recur. First, every flow measured in final goods is grossed up by a dimensionless factor $1\pm\zeta\Omega^j$, because the activity carries embodied material along with the resources it absorbs: consumption in~\eqref{eq:sp:sum:C}, extraction and exploration effort in~\eqref{eq:sp:sum:N} and~\eqref{eq:sp:sum:D}, treatment in~\eqref{eq:sp:sum:varpi}, and output itself wherever $F_K$ or $F_R$ appears. These wedges are what the composition of activities/uses does to the allocation, and they vanish with $\zeta$. Second, the ledger charges each tonne entering $R$ exactly once, whatever route it took: this is why extraction pays $p^W$ per tonne in~\eqref{eq:sp:sum:N} while exploration pays only for the mass it displaces in~\eqref{eq:sp:sum:D}, and why the recovered tonne in~\eqref{eq:sp:sum:varpi} is worth $\Psi-p^W$ rather than $\Psi$.

The system closes in the usual way. Given the six states and the six costates $\left(q,p^S,p^X,p^P,p^M,p^{\mathcal W}\right)$, the eight first-order conditions~\eqref{eq:sp:sum:Omega}--\eqref{eq:sp:sum:KR} together with the three static constraints of Problem~\ref{prob:sp} form eleven equations in eleven unknowns: the eight controls and the three prices $\lambda$, $p^W$ and $\zeta$ that carry no state of their own --- with the waste condition~\eqref{eq:sp:sum:W} now delivering $p^W$ at once as $-p^{\mathcal W}$, so that only $\lambda$ and $\zeta$ are genuinely static. Four of the eleven are equations only branch by branch: on the extraction, exploration, treatment and recycling-capital margins the system selects between a corner and an interior root according to~\eqref{eq:sp:sum:Nsys}, \eqref{eq:sp:sum:Dsys}, \eqref{eq:sp:sum:varpi} and~\eqref{eq:sp:sum:KR}, which is one equation either way and is what makes the count work in every phase. The costate recursions~\eqref{eq:sp:sum:K}--\eqref{eq:sp:sum:Wstock} and the transitions~\eqref{eq:sp:states-motion} then propagate the twelve remaining objects, with the consumption rule~\eqref{eq:sp:sum:Cgrowth} following from~\eqref{eq:sp:sum:C} and the definition of $r$ rather than adding a restriction of its own. The timing is recursive within each period: the handled flow $H_t=\mu\mathcal W_t$ is predetermined, so the material block never contains a within-period fixed point --- period-$t$ waste feeds period-$t{+}1$ recycling, never its own.
